\subsection{FUNCTIONAL RELATIONS}

Let R be an assembly and x a letter. Let y and z be distinct letters which are distinct from x and which do not appear in R. Let \(y'\) , \(z'\) be two other letters with the same properties. By CS8, CS9 (§4, no. 1), CS2, CS5 (§1, no. 2), and CS6 (§3, no. 4), the assemblies

\[
(\forall y) (\forall z) (((y | x) R \text { and } (z | x) R) \Longrightarrow (y = z))
\]

and

\[
(\forall y ^ {\prime}) (\forall z ^ {\prime}) (((y ^ {\prime} | x) R \text { and } (z ^ {\prime} | x) R) \Longrightarrow (y ^ {\prime} = z ^ {\prime}))
\]

are identical. If R is a relation in G, the assembly thus defined is a relation in G which is denoted by ``there exists at most one x such that R''; the letter x does not appear in this relation. When this relation is a theorem in G, R is said to be single-valued in X in G. To prove that R is single-valued in the theory G, it is enough to prove y = z in the theory obtained by adjoining to G the axioms \((y|x)R\) and \((z|x)R\) , where y and z are distinct letters which are distinct from x and appear neither in R nor in the explicit axioms of G.

C45. Let R be a relation in C, and let x be a letter which is not a constant of C. If R is single-valued in x in C, then \(\mathbf{R} \Longrightarrow (\mathbf{x} = \tau_{\mathbf{x}}(\mathbf{R}))\) is a theorem in C. Conversely if, for some term T in C which does not contain x, \(\mathbf{R} \Longrightarrow (\mathbf{x} = \mathbf{T})\) is a theorem in C, then R is single-valued in x in C.

Suppose R is single-valued in x in T, and let us show that

\[
R \Longrightarrow (x = \tau_ {x} (R))
\]

is a theorem in \(\mathcal{C}\) . Adjoin the hypothesis \(R\) . Then \((\tau_{x}(R)|x)R\) is true by S5, and hence `` \(R\) and \((\tau_{x}(R)|x)R\) '' is true. Now, since \(R\) is single-valued in \(x\) ,

\[
(R \text {   and   } (\tau_ {x} (R) | x) R) \Longrightarrow (x = \tau_ {x} (R))
\]

is a theorem in \(\mathfrak{G}\) by C30 (§4, no. 3). Therefore \(x = \tau_x(R)\) is true. ¶ Conversely, suppose that \(R \Rightarrow (x = T)\) is a theorem in \(\mathfrak{G}\) . Let \(y, z\) be distinct letters which are distinct from \(x\) and which appear neither in \(R\) nor in the explicit axioms of \(\mathfrak{G}\) . Since \(x\) is not a constant of \(\mathfrak{G}\) and does not appear in \(T\) , the relations

\[
(y | x) R \Longrightarrow (y = T), \quad (z | x) R \Longrightarrow (z = T)
\]

are theorems in \(\mathfrak{G}\) . Adjoin the hypotheses \((y|x)R\) and \((z|x)R\) . Then \(y = T\) and \(z = T\) are true, hence \(y = z\) is true.

Let R be a relation in C. The relation

``\((\exists x)R\) and there exists at most one x such that R''

is denoted by ``there exists exactly one x such that R''. If this relation is a theorem in C, R is said to be a functional relation in x in the theory C.

C46. Let R be a relation in G, and let x be a letter which is not a constant of G. If R is functional in x in G, then \(\mathbf{R} \Longleftrightarrow (\mathbf{x} = \tau_{\mathbf{x}}(\mathbf{R}))\) is a theorem in G. Conversely, if for some term T in G which does not contain x,

\[
\boldsymbol {R} \Longleftrightarrow (\boldsymbol {X} = \boldsymbol {T})
\]

is a theorem in \(\mathfrak{G}\) , then \(\pmb{R}\) is functional in \(x\) in \(\mathfrak{G}\) .

Suppose R is functional in x in C. Then \(\boldsymbol{R} \Longrightarrow (\boldsymbol{x} = \tau_{\boldsymbol{x}}(\boldsymbol{R}))\) is a theorem in C by C45. On the other hand, \((\exists \boldsymbol{x})\boldsymbol{R}\) is a theorem in T. By S6 the relation

\[
(x = \tau_ {x} (R)) \Longrightarrow (R \Longleftrightarrow (\exists x) R)
\]

is a theorem in \(\mathfrak{C}\) . If we adjoin the hypothesis \(x = \tau_x(R)\) , it follows that \(R\) is true. Therefore \((x = \tau_x(R)) \Longrightarrow R\) is a theorem in \(\mathfrak{C}\) . Conversely, if \(R \Longleftrightarrow (x = T)\) is a theorem in \(\mathfrak{C}\) , then \(R\) is single-valued in \(x\) in \(\mathfrak{C}\) , by C45. Moreover, \((T|x)R \Longleftrightarrow (T = T)\) is a theorem in \(\mathfrak{C}\) ; hence \((T|x)R\) and therefore \((\exists x)R\) are theorems in \(\mathfrak{C}\) .

If a relation R is functional in x in T, then R is equivalent to the relation \(\boldsymbol{x} = \tau_{\boldsymbol{x}}(\boldsymbol{R})\) , which is often more manageable. Generally an abbreviating symbol \(\Sigma\) is introduced to represent the term \(\tau_{\boldsymbol{x}}(\boldsymbol{R})\) . Such a symbol is called a functional symbol in G.

Intuitively, \(\Sigma\) represents the unique object which has the property defined by \(R\) . * For example, in a theory where `` \(y\) is a real number \(\geqslant 0\) '' is a theorem, the relation `` \(x\) is a real number \(\geqslant 0\) and \(y = x^{2}\) '' is functional in \(x\) . The corresponding functional symbol is taken to be either \(\sqrt{y}\) or \(y^{1/2}\) . *

C47. Let x be a letter which is not a constant of \(\mathfrak{G}\) , and let \(R\{x\}\) and \(S\{x\}\) be two relations in \(\mathfrak{G}\) . If \(R\{x\}\) is functional in x in \(\mathfrak{G}\) , then the relation \(S\{\tau_{x}(R)\}\) is equivalent to \((\exists x)(R\{x\}\) and \(S\{x\})\) .

For it follows from C46 and C43 that \((\pmb{R}\{\pmb{x}\}\) and \(S\{\pmb{x}\})\) is equivalent to \((\pmb{R}\{\pmb{x}\}\) and \(S\{\tau_{\pmb{x}}(\pmb{R})\}\) ; since \(S\{\tau_{\pmb{x}}(\pmb{R})\}\) does not contain \(\pmb{x}\) ,

\[
(\exists x) (R \{x \} \text {   and   } S \{\tau_ {x} (R) \})
\]

is equivalent to \((S\{\tau_x(R)\}\) and \((\exists x)R)\) by C33 (§ 4, no. 3); and the result follows from the fact that \((\exists x)R\) is true, because \(R\) is functional in \(x\) .

\clearpage
